\section{Exact ordinary moment laws for the full theta family}
\label{gm-section}

Integral Taylor moments alone do not identify translation by
$p$-power torsion with multiplication of a measure by a character.
We prove that identification for the actual theta functions.  A
finite expansion into elliptic theta functions supplies the
identity; the independent normalization already proved supplies
the bound on the measure.  No root of a translated identity is
taken.

\subsection{An orthogonal frame with prime-to-$p$ index}

Retain $k=\mathbb Q(\sqrt{-D})$, $D\equiv7\pmod8$, and the
fixed abstract isomorphism $\iota:\mathbb C\simeq\mathbb C_p$.
Assume $p\ge5$ splits in $k$.  Let $V_k$ be a positive definite
Hermitian $k$-space of dimension $g$, with form $h$ linear in its
first variable, and let $\Lambda$ be an integral unimodular
$\mathcal O_k$-lattice.  On its complexification use
\[
 H=\frac{2}{\sqrt D}h,\quad E=\operatorname{Im}H,\quad
 Q(x)=h(x,x),\quad \alpha(x)=e^{\pi iQ(x)},\quad
 W_x=\alpha(x)U_x.                                      \tag{GM1}
\]
The lattice is self-dual for $E$, and $\alpha$ restricts to
its norm semicharacter.  The trace lattices (TE2) have precisely
this description after choosing their rational $k$-structure;
their common complex scaling is included in that identification.

\begin{lemma}\label{gm-orthogonal-frame}
There are $x_1,\ldots,x_g\in\Lambda$ such that
\[
 h(x_i,x_j)=a_i\delta_{ij},\qquad
 a_i\in\mathbb Z_{>0},\quad p\nmid a_i.                 \tag{GM2}
\]
For $\Lambda_0=\bigoplus_i\mathcal O_kx_i$,
the finite index $[\Lambda:\Lambda_0]$ is prime to $p$,
and
\[
 \Lambda_0^\#=\bigoplus_i a_i^{-1}\mathcal O_kx_i,
                  \qquad \Lambda_0\subset\Lambda\subset\Lambda_0^\#.
                                                               \tag{GM3}
\]
Here the dual can be taken for $h$ or for $E$.
\end{lemma}
\begin{proof}
At $p$ the ring $\mathcal O_k\otimes\mathbb Z_p$ is
$\mathbb Z_p\times\mathbb Z_p$, with the factors interchanged
by conjugation.  A unimodular Hermitian pairing is a perfect
pairing of its two factors.  A basis of one factor and its dual
basis of the other give an orthonormal basis of
$\Lambda\otimes\mathbb Z_p$.

Approximate that basis sufficiently closely by vectors in
$\Lambda$, which is dense in its $p$-adic completion.  Apply
Hermitian Gram--Schmidt over $k$.  Every successive denominator
is a $p$-adic unit, because the approximating Gram matrix is
congruent to $1$ modulo $p$.  The resulting orthogonal vectors
belong to $\Lambda\otimes\mathbb Z_{(p)}$ and still form a
basis at $p$.  Multiply each by a positive integer prime to $p$
to put it in $\Lambda$.  Its norm is then in
$\mathcal O_k\cap\mathbb Q=\mathbb Z$, is positive by
definiteness, and is a $p$-adic unit.  This proves (GM2).
The local basis property implies that the finite quotient
$\Lambda/\Lambda_0$ has order prime to $p$.
The displayed formula for the dual follows from the diagonal
Gram matrix.  Integrality of $h$ on $\Lambda$ gives the second
inclusion in (GM3).  Hermitian and alternating duality agree
because the inverse different of $k/\mathbb Q$ is
$(\sqrt{-D})^{-1}\mathcal O_k$.
\end{proof}

Let $\theta$ be the normalized even elliptic theta function of
Section~\ref{ct-section}, for the lattice $\mathcal O_k$.
Use its single character-normalized ordinary period $\omega$.
Let $\mathcal T(\Lambda)$ denote the entire functions satisfying
$W_lF=F$ for every $l\in\Lambda$.

\begin{proposition}\label{gm-finite-expansion}
For every $F\in\mathcal T(\Lambda)$ there are finitely many
complex coefficients $c_j$ and $t_{ji}\in k$ such that
$a_it_{ji}\in\mathcal O_k$ and
\[
 F\left(\sum_i z_ix_i\right)
       =\sum_j c_j\prod_{i=1}^g(U_{t_{ji}}\theta)(z_i)^{a_i}.
                                                               \tag{GM4}
\]
The $U$ in each factor is the elliptic operator, with Hermitian
form $2z\bar w/\sqrt D$.  Thus all its translates have order
prime to $p$.  The coefficients need not be assumed integral.
\end{proposition}
\begin{proof}
Put $G(\sum_i z_ix_i)=\prod_i\theta(z_i)^{a_i}$.
It is nonzero, belongs to $\mathcal T(\Lambda_0)$, and
$G(0)=1$.  The commutator pairing
\[
 (\Lambda/\Lambda_0)\times(\Lambda_0^\#/\Lambda)
       \longrightarrow\mathbb C^\times,
             (l,w)\longmapsto e^{-2\pi iE(l,w)}              \tag{GM5}
\]
is perfect.  Indeed, vanishing of the pairing with all
$\Lambda_0^\#$ means membership in its dual $\Lambda_0$;
vanishing with all $\Lambda$ means membership in
$\Lambda^\#=\Lambda$.  The two finite groups have the
same order by their dual indices.

Choose representatives $l$ and $w$ of these groups, with
the zero coset represented by zero, and set
\[
                      G_w=\sum_{l\in\Lambda/\Lambda_0}W_lU_wG.
                                                               \tag{GM6}
\]
Since $w\in\Lambda_0^\#$, $U_wG$ has the same
$\Lambda_0$ multiplier as $G$.  The operators $W_l$ on
$\Lambda$ are additive, so (GM6) is well-defined and belongs
to $\mathcal T(\Lambda)$.  Commuting $U_{-w}$ past $W_l$
and using (GM5) gives
\[
 \sum_{w\in\Lambda_0^\#/\Lambda}U_{-w}G_w
    =\sum_l\left(\sum_w e^{-2\pi iE(l,w)}\right)W_lG
    =|\Lambda_0^\#/\Lambda|G\ne0.                            \tag{GM7}
\]
Hence some $G_w$ is nonzero.  By
Lemma~\ref{tc-covolume}, $\mathcal T(\Lambda)$ has dimension
one, so $F=cG_w$ for a scalar $c$ (including $c=0$).
For each summand in (GM6), the Heisenberg composition rule gives
\[
 W_lU_wG=\alpha(l)e^{-\pi iE(l,w)}U_{l+w}G.
\]
The coordinates of $l+w$ belong to $a_i^{-1}\mathcal O_k$
by (GM3).  Orthogonality of the frame identifies its translate
with the product in (GM4).  Absorb the displayed scalar into
$c_j$.  The phases are roots of unity, since $E(l,w)$ is
rational.  This proves the assertion.
For $g=0$ the empty product is $1$ and the assertion is immediate.
\end{proof}

\subsection{The measure and its torsion characters}

Write $\iota_-:k\to\mathbb Q_p$ for $b\mapsto\iota(\bar b)$,
and put
\[
                 X_\Lambda=\Lambda\otimes_{\mathcal O_k,\iota_-}
                                            \mathbb Z_p.        \tag{GM8}
\]
This is the conjugate ordinary component of the lattice, a free
$\mathbb Z_p$-module of rank $g$.  For $v\in\Lambda$ the
function
\[
              L_v(z)=\iota(h(v,z)),\qquad z\in\Lambda,
                                                               \tag{GM9}
\]
extends to a $\mathbb Z_p$-linear function on $X_\Lambda$:
in the frame (GM2), it is
$\sum_i a_i\iota(v_i)z_i$, where $z_i$ denotes the
$\iota_-$-coordinate.  This convention is essential; the
Hermitian adjoint in the point moments uses this conjugate
component.

For $y=S\kappa l/p^n$, with $l\in\Lambda$,
$S\in\mathbb Z_{>0}$ prime to $p$, and
\[
        \kappa\in\mathcal O_k,
                   \iota(\bar\kappa)\in p^{2n}\mathbb Z_p,
                                                               \tag{GM10}
\]
define on rational lattice representatives
\[
                  \chi_y(z)=\iota\bigl(e^{-2\pi iE(z,y)}\bigr).
                                                               \tag{GM11}
\]
It induces a character of $X_\Lambda/p^nX_\Lambda$.
For example, in an orthogonal frame in which all $y_i$ have
the strong divisibility (CT24),
\[
        \chi_y(z)=\prod_i
          \iota\bigl(e^{-2\pi iE_1(1,y_i)}\bigr)^{a_i z_i\bmod p^n},
             \qquad E_1=\operatorname{Im}(2z\bar w/\sqrt D).
                                                               \tag{GM12}
\]
To verify this formula, adjointness gives
$E_1(b,y_i)=E_1(1,\bar b y_i)$.
On the formal $p^n$-torsion, multiplication by $\bar b$ is
multiplication by $\iota(\bar b)$ modulo $p^n$; this is
Proposition~\ref{ct-character-normalization}.  It also applies
to $b\in\mathcal O_k\otimes\mathbb Z_{(p)}$, by inverting
an integer prime to $p$.  Equation (GM12) follows and proves
the assertion about factorization.

\begin{theorem}\label{gm-moment-law}
Let $F\in\mathcal T(\Lambda)$ and let $t\in V_k$ have
order prime to $p$ modulo $\Lambda$.
There is a unique bounded $\mathbb C_p$-valued measure
$\mu_{F,t}$ on $X_\Lambda$ such that, for all
$v_1,\ldots,v_r\in\Lambda$,
\[
 \widehat{U_tF}\left(\omega\sum_jv_j\log(1+T_j)\right)
     =\int_{X_\Lambda}\prod_j(1+T_j)^{L_{v_j}(z)}
                                                  \,d\mu_{F,t}(z).
                                                               \tag{GM13}
\]
There is an integer $S>0$ prime to $p$, independent of $n$,
$l$ and $\kappa$, such that for every $y$ as in (GM10),
\[
 \widehat{W_yU_tF}\left(\omega\sum_jv_j\log(1+T_j)\right)
     =\int_{X_\Lambda}\chi_y(z)
                \prod_j(1+T_j)^{L_{v_j}(z)}\,d\mu_{F,t}(z).
                                                               \tag{GM14}
\]
In particular the same measure supplies all derivatives and
all these torsion characters.  For each fixed $n$, the identity
also holds for the additive subgroup generated by these $y$.
If the period-scaled moments of
$U_tF$ are integral on $\Lambda$, then $\|\mu_{F,t}\|\le1$.
\end{theorem}
\begin{proof}
Use (GM4).  The coordinates $t_i$ of $t$ have prime-to-$p$
denominators.  In these coordinates $U_t$ replaces $t_{ji}$
by $t_i+t_{ji}$ and introduces only a root of unity into each
coefficient.  Every transported elliptic series is bounded
and integral by Proposition~\ref{ct-translates-integral}.
Their powers, products and the finite sum are bounded, by
$\max_j|\iota(c_j)|_p$.  By the multivariable moment theorem
of Section~\ref{ot-section}, this series is the Amice
transform of a bounded measure on $\mathbb Z_p^g$.
In the frame, $L_{x_i}(z)=a_i z_i$.  Since $a_i$ is a
$p$-adic unit, push this measure forward by
$(A_i)\mapsto(A_i/a_i)$ to obtain $\mu_{F,t}$.

For general directions $v_j$, the coefficients $\iota(v_{ji})$
lie in $\mathbb Z_p$.  Substitution (LM3) proves (GM13).
The frame directions already determine the measure uniquely,
because the binomial polynomials are dense in the continuous
functions on $\mathbb Z_p^g$, as proved in
Section~\ref{ot-section}.

Choose $s>0$ prime to $p$ so that every
$s(t_i+t_{ji})$ is in $\mathcal O_k$.  Choose $d>0$ prime
to $p$ so that $d\Lambda\subset\Lambda_0$, and put $S=sd$.
If $l=\sum l_ix_i$, then
\[
          y_i=\frac{s\nu_i}{p^n},\qquad
          \nu_i=d\kappa l_i\in\mathcal O_k,
                      \iota(\bar\nu_i)\in p^{2n}\mathbb Z_p.
                                                               \tag{GM15}
\]
Theorem~\ref{ct-power-translation} applies to every elliptic
factor with this same $s$.  Multiplying its identities,
taking the powers $a_i$, and adding the finitely many terms
proves the multivariable substitution identity
\[
      \widehat{W_yU_tF}(T_1,\ldots,T_g)
          =\widehat{U_tF}\bigl(\zeta_1(1+T_1)-1,
                         \ldots,\zeta_g(1+T_g)-1\bigr),
 \quad \zeta_i=\iota(e^{-2\pi iE_1(1,y_i)}).              \tag{GM16}
\]
In (GM16), the hats stand for the frame expansion
$\omega\sum_i x_i\log(1+T_i)$.
The Weyl phase is exactly the product of the elliptic phases
raised to $a_i$, since $Q(y)=\sum_i a_i|y_i|^2$.
Thus no root of unity was discarded in deriving (GM16).
The convergent substitutions (CT21) commute with finite
products and sums.  Under the Amice transform they multiply
the measure by $\prod_i\zeta_i^{A_i}$, which is precisely
(GM12) after $A_i=a_i z_i$.  This proves (GM14) for frame
directions, and (LM3) proves it for all directions.
For a finite sum or negative of the permitted lifts at one
fixed level $n$, each
coordinate in (GM15) still has an integral numerator with the
same conjugate divisibility.  Apply the elliptic law directly
to these summed numerators; the same proof gives (GM14) on
the indicated additive subgroup.
Finally, integrality on the frame directions says that the
Amice transform has coefficient norm at most one.  The norm
identity in the moment theorem and the invertible coordinate
change $A_i=a_i z_i$ give $\|\mu_{F,t}\|\le1$.
\end{proof}

The characters appearing here include all characters of
$X_\Lambda/p^nX_\Lambda$.  Indeed choose $\kappa$ with
$\iota(\kappa)\equiv1\pmod{p^{2n}}$ and
$\iota(\bar\kappa)\equiv0\pmod{p^{2n}}$ by the Chinese
remainder theorem.  The alternating pairing of a self-dual
lattice pairs its two ordinary components perfectly modulo
$p^n$.  In (GM2) this is the product of the elementary elliptic
pairings, multiplied by the units $a_i$.  Multiplication by
$S$ is invertible modulo $p^n$, and $l\in\Lambda$ surjects
onto either component modulo $p^n$.  Formula (GM11) therefore
runs through every character.

\subsection{The bound for the normalized arithmetic family}

\begin{theorem}\label{gm-arithmetic-unit-bound}
Retain the hypotheses and the even dimensions of
Section~\ref{te-section}.  Let $p\ge5$ split in $k$, and
assume $p$ is prime to the elliptic ideal norm $a$.
For all unimodular $\mathcal O_E$-lattices $L,M$, the
normalized arithmetic vector $e(L,M)$ of (TE15)
and every translate of order prime to $p$ have the exact
laws (GM13)--(GM14), with $\|\mu_{e(L,M),t}\|\le1$.
The period is the common period $\omega$ of
Proposition~\ref{ci-common-period}; it is not selected again
after choosing $L,M$, a derivative, or a torsion point.
\end{theorem}
\begin{proof}
First suppose $L$ and $M$ are standard at every prime over
$p$.  The base vector is the tensor product of the normalized
trace-lattice theta vector.  Theorem~\ref{ts-integral-theta}
and multiplication of Amice transforms give its integral
moments and those of every prime-to-$p$ translate.
Each of the two finite averaging operators defining $e(L,M)$
has index prime to $p$, because its lattices have equal
$p$-adic completions.  Proposition~\ref{lm-finite-intertwiner}
preserves the asserted integrality.

The determinant characters in its definition are $p$-adic
units.  Here is the exact argument for this character convention.
For a norm-one fractional ideal $J$, choose $b>0$ with
$J^b=(z)$, using finiteness of the ideal class group.  Then
$J^{2b}=(z/\bar z)$, with $z/\bar z$ of norm one.  The
principal-ideal convention (TE16) gives
\[
            \chi(J)^{2b}
                   =\prod_{\eta\in\Sigma}\eta(z/\bar z)^{-1}.
                                                               \tag{GM17}
\]
If $J$ is a unit above $p$, both $z$ and $\bar z$ are
units there.  Every factor on the right is a $p$-adic unit
at every transported place above $p$.  Taking valuations in
(GM17) proves the assertion.  The relative determinant ideals
of the standard local lattices have this property.

For arbitrary $L,M$, choose rational unitary transformations
$\gamma,\gamma'$ such that $\gamma L$ and $\gamma'M$
are standard at every prime over $p$.  This uses only weak
approximation: Proposition~\ref{acl-local-basis} gives a local
unitary transformation at each of those places.  Perturb it
within the open set preserving the specified lattice so that
$1+g$ is invertible.  Such a perturbation exists by multiplying
by local norm-one scalars near $1$ and avoiding the finitely
many roots of $\det(1+zg)$.  Its Cayley transform
$(1-g)(1+g)^{-1}$ is skew-Hermitian.  Weak approximation in
the finite-dimensional $F$-space of skew-Hermitian matrices,
followed by the inverse Cayley transform, gives the required
rational transformations simultaneously.  Nonvanishing of
the denominator and the prescribed lattice images are open
conditions, so they survive sufficiently close approximation.

The exact covariance proved in Section~\ref{te-section}
identifies $e(\gamma L,\gamma'M)$ with the pullback of
$e(L,M)$ by $z\mapsto\gamma'^*z\gamma$.  This unitary
map carries the corresponding period lattice onto the original
one.  Thus the integral moments just proved, including all
prime-to-$p$ translates, pull back to those on $\Lambda(L,M)$.
Apply Theorem~\ref{gm-moment-law} and its norm conclusion.
The equalities of periods under ideal isogenies and common
scaling were proved in Section~\ref{ci-section}; no assertion
about the $p$-adic absolute value of the complex scaling is
needed.
\end{proof}

\begin{corollary}\label{gm-algebraic-jets}
Let $\delta_0$ be the discriminant twelfth root for the
principal elliptic lattice.  For the family in
Theorem~\ref{gm-arithmetic-unit-bound}, every rational-lattice
translate $t\in\mathbb Q\Lambda(L,M)$ has algebraic jets
\[
       \delta_0^{-k}D_{v_1}\cdots D_{v_k}
                           (U_te(L,M))(0)\in\overline{\mathbb Q}
       \qquad(v_i\in\mathbb Q\Lambda(L,M)).                 \tag{GM18}
\]
The assertion also permits algebraic linear combinations of
these directions.
\end{corollary}
\begin{proof}
Fix an embedding of $\overline{\mathbb Q}$ into
$\mathbb C_p$, and let $\iota$ run through all its extensions
to an abstract isomorphism $\mathbb C\simeq\mathbb C_p$.
The previous proof applies to each one.  The ordinary period
can change with $\iota$, but always has the form
$\omega=\Omega/\iota(\delta_0)$ with $|\Omega|_p=1$.
Choose a sufficiently large split prime $p$ avoiding the
denominators of $t,v_1,\ldots,v_k$ and the elliptic ideal
norm.  Clearing those denominators, the moment formula and
the unit bound give an upper bound independent of $\iota$
for the absolute value of the transport of (GM18).  In fact
the bound is one if the directions belong to the lattice.

We recall why this forces algebraicity.  If a complex number
$c$ were transcendental, fix one such isomorphism $\iota_0$
and choose a transcendental $a\in\mathbb C_p$, over the
image of $\overline{\mathbb Q}$, with arbitrarily large
absolute value.  Such an $a$ exists because that algebraic
closure is countable whereas every nonempty open ball is
uncountable.  The isomorphism
$\overline{\mathbb Q}(c)\to
\overline{\mathbb Q}(\iota_0^{-1}a)$ sending $c$ to
$\iota_0^{-1}a$ extends to an automorphism of $\mathbb C$:
extend transcendence bases, then extend over their algebraic
closures.  Composing with $\iota_0$ contradicts the uniform
upper bound.  Thus the number in (GM18) is algebraic.
Multilinearity gives the final assertion.

Only existence of a split prime outside a finite list is
used here.  It also has an elementary proof in this quadratic
case: choose $m$ divisible by $2D$ and all primes in that
list.  Every prime divisor $q$ of $Dm^2+1>1$ is outside
the list, is prime to $2D$, and satisfies
$(Dm)^2\equiv-D\pmod q$.  Hence $q$ splits in $k$.
\end{proof}

This theorem supplies the character identities needed for finite
Fourier restriction of the actual measures.  It does not yet
assert that a particular local test function selects the two
matrix supports in (CV8); that identification also requires the
explicit finite translation operator defining the arithmetic
column.


